\section{Ordinary trace volumes on a nonclosed ruled surface}
\label{sec:nonclosed-ruled-trace-masses}

The preceding product formula computes the mass of a separated representative. When the background itself varies along the curve, even bounded changes of the representative can change its mass. On a product of projective lines we can optimize over the entire ordinary singularity type, including potentials depending on both factors. The potential will be a logarithmic pole along a divisor plus an unbounded residual with zero Lelong number. Its multiplier ideals see only the imposed divisor, whereas its ordinary type also retains the residual pole.

In this section let $Y=\mathbb P^1\times\mathbb P^1$, let $w$ be the probability Fubini--Study form on $\mathbb P^1$, and denote its pullbacks by $w_1,w_2$. Fix $p\in\mathbb P^1$, and let $G_p\leq0$ be the Green potential satisfying $\ddc G_p=[p]-w$ and $\sup G_p=0$. Functions on the first factor are also viewed as functions on $Y$. Put
\[
 D=\{p\}\times\mathbb P^1,\qquad
 U=1-(1-G_p)^\alpha,\qquad r=bG_p+U,
 \qquad 0<\alpha<1,
\]
where $b\geq0$. The term $bG_p$ is identically zero when $b=0$. Choose $c>0$ and a smooth positive function $a$ on the first curve, and set
\begin{equation}\label{eq:ruled-trace-background}
 \theta=cw_1+aw_2.
\end{equation}
Thus $\theta$ is Hermitian and is nonclosed precisely when $a$ is nonconstant. The function $r$ specifies an ordinary type; whether that type has $\theta$-plurisubharmonic representatives is decided in \cref{thm:nonclosed-ruled-ordinary-volumes}.

\begin{lemma}\label{lem:ruled-sublog-small-curvature}
There are $U_L\sim U$ with $\sup U_L=0$, $\ddc U_L\geq-\epsilon_Lw$ and $\epsilon_L\downarrow0$, such that $U_L\to0$ in $L^1(\mathbb P^1)$. Each has zero Lelong number at $p$ and is locally bounded off $p$.
\end{lemma}
\begin{proof}
For $t\leq0$ put $\chi(t)=1-(1-t)^\alpha$. Then
\[
 \chi'(t)=\alpha(1-t)^{\alpha-1}>0,\qquad
 \chi''(t)=\alpha(1-\alpha)(1-t)^{\alpha-2}>0.
\]
For $L>0$ set $\epsilon_L=\chi'(-L)$, $C_L=-\epsilon_LL-\chi(-L)$, and define
\[
 \chi_L(t)=
 \begin{cases}
  \chi(t)+C_L,&t<-L,\\
  \epsilon_Lt,&-L\leq t\leq0.
 \end{cases}
\]
The values and first derivatives match, and $\chi_L$ is convex and increasing with derivative at most $\epsilon_L$. Off $p$, the chain rule gives $\ddc\chi_L(G_p)\geq-\epsilon_Lw$. Canonical extension across $p$ retains this bound by the removable-singularity theorem, so $U_L=\chi_L(G_p)$ has the required curvature. Near $p$ it differs from $U$ by $C_L$, and elsewhere both are bounded. Its growth is sublogarithmic, hence its Lelong number is zero.

Writing $x=1+L$, we have $C_L=(1-\alpha)x^\alpha+\alpha x^{\alpha-1}-1\geq0$. Also $\chi(t)\leq\epsilon_Lt$ for $-L\leq t\leq0$, by integrating $\chi'\geq\epsilon_L$ on this interval. It follows that $U\leq U_L\leq0$. Since $\epsilon_L\to0$, the functions tend to zero off $p$, and dominated convergence gives the $L^1$ limit.
\end{proof}

\begin{theorem}\label{thm:nonclosed-ruled-ordinary-volumes}\index{Trace!nonclosed ruled surface}\index{Singularity volume!unbounded ruled type}
The ordinary type of $r$ has $\theta$-plurisubharmonic representatives if and only if $b<c$. In this range it contains a strict current: some $v\sim r$ satisfies $\theta+\ddc v\geq\delta(w_1+w_2)$ with $\delta>0$. Every $v\in\PSH(Y,\theta)$ with $v\sim r$ has positive finite extended mass, given by
\begin{equation}\label{eq:nonclosed-ruled-affine-mass}
 \int_Y\mathrm{MA}^{\mathrm{ext}}_\theta(v)
   =\int_Y\theta^2-2b a(p)+2\int_Yv\,\ddc\theta.
\end{equation}
The volumes of \cref{def:singularity-volumes} are
\begin{equation}\label{eq:nonclosed-ruled-ordinary-volumes}
 \underline{\operatorname{vol}}_\theta(v)=2(c-b)\min a,
 \qquad
 \overline{\operatorname{vol}}_\theta(v)=2(c-b)\max a.
\end{equation}
Both values are limits of masses of representatives of this ordinary type.
\end{theorem}
\begin{proof}
We first remove the logarithmic divisor. If $v\sim r$, then $\nu(v,D)=b$, so \cref{lem:nonclosed-divisorial-removal} gives $\theta+\ddc v-b[D]\geq0$. Taking $q=v-bG_p$ in $L^1$ and then its quasi-plurisubharmonic representative gives
\[
 q\in\PSH(Y,\beta),\qquad q\sim U,\qquad
 \beta=kw_1+aw_2,\qquad k=c-b.
\]
Indeed the ordinary inequalities hold off $D$ and extend canonically across it. Conversely such a $q$ gives $v=bG_p+q$ with
\[
 \theta+\ddc v=b[D]+\beta+\ddc q.
\]
Their locally bounded products agree off $D$, and both extended measures give $D$ zero mass. Thus their extended measures agree everywhere.

If $k>0$, choose $L$ with $\epsilon_L<k/2$. The potential $bG_p+U_L$ belongs to the specified type and satisfies
\[
 \theta+\ddc(bG_p+U_L)
 \geq b[D]+(k-\epsilon_L)w_1+aw_2
 \geq\min(k/2,\min a)(w_1+w_2).
\]
For necessity, average $q$ over the unitary group acting transitively on the second projective line, with its Haar probability measure. The background $\beta$ is invariant, so the average is $\beta$-plurisubharmonic and invariant, hence a function $q_0$ on the first factor, with $q_0\sim U$ and $kw+\ddc q_0\geq0$. One can first average in $L^1$ and then take the canonical representative; the curvature inequality and the ordinary bounds pass to it. The positive curve measure has total mass $k$, excluding $k<0$. When $k=0$, $q_0$ is harmonic on the compact curve and therefore constant, contradicting $q_0\sim U$.

It remains to compute the mass for every competitor, before using the average. Choose a closed K\"ahler form $K=C(w_1+w_2)\geq\beta$ and a fixed $u_0=U_L$ as above. The closed positive current $S_0=K+\ddc u_0$ has no coefficient mass on $D$, since its first-factor curvature has no atom at $p$. The product formula of \cref{prop:curve-product-extended-ma} gives
\[
 \langle S_0\rangle=S_0,\qquad
 \langle S_0^2\rangle=2C(Cw_1+\ddc u_0)\wedge w_2,
 \qquad \int_Y\langle S_0^2\rangle=\int_YK^2.
\]
For $S=K+\ddc q$, both $S$ and $S_0$ are closed positive currents in the same K\"ahler class, with bounded difference of potentials and small unbounded locus $D$. The classical closed-class comparison of non-pluripolar products, applied in both singularity orders in degrees one and two, gives
\[
 \{\langle S\rangle\}=\{K\},\qquad
 \int_Y\langle S^2\rangle=\int_YK^2.
\]
Here the degree-one non-pluripolar current is the restriction of $S$ off $D$ and is closed. Hence the positive current $S-\langle S\rangle$ has zero trace against $K$, and vanishes. This proves $\langle S\rangle=S$. Alternatively, absence of polar coefficient mass follows from \cref{prop:bounded-type-polar-invariance}. The closed comparison is used only for $K$.

Put $R=\beta-K$. Bedford--Taylor expansion off $D$, followed by extension by zero, now yields
\[
 \mathrm{MA}^{\mathrm{ext}}_\beta(q)
     =\langle S^2\rangle+2R\wedge S+R^2.
\]
The first term is finite by closed comparison, the second has finite variation because $S$ has finite trace, and the last is smooth. Since none has coefficient mass on $D$, this is also an identity on $Y$. Distributional Stokes therefore gives
\begin{equation}\label{eq:ruled-residual-rough-mass}
 \int_Y\mathrm{MA}^{\mathrm{ext}}_\beta(q)
 =\int_Y\beta^2+2\int_Yq\,\ddc\beta
 =2\int_{\mathbb P^1}a(kw+\ddc q_0).
\end{equation}
The pairing uses only $q\in L^1$. Averaging preserves it, since $\ddc\beta=\ddc a\wedge w_2$. The last measure is positive of total mass $k$, proving the bounds in \eqref{eq:nonclosed-ruled-ordinary-volumes} for all original competitors. Substituting $q=v-bG_p$ and $\int G_p\ddc a=a(p)-\int aw$ proves \eqref{eq:nonclosed-ruled-affine-mass}.

To approach either bound, choose a point $x$ where $a$ takes the corresponding extreme value, smooth strictly positive measures $\zeta_j$ of mass $k$ tending weakly to $k\delta_x$, and smooth $g_j$ with $\ddc g_j=\zeta_j-kw$. Choose $L_j$ so that $\epsilon_{L_j}$ is less than half the minimum density of $\zeta_j$ relative to $w$ and $|\int U_{L_j}\ddc a|<1/j$. Then
\[
 v_j=bG_p+g_j+U_{L_j}\sim r,\qquad
 \beta+\ddc(g_j+U_{L_j})\geq\tfrac12\zeta_j+aw_2>0.
\]
Their masses tend to $2ka(x)$ by \eqref{eq:ruled-residual-rough-mass}.
\end{proof}

No finite first moment is used here. For the reference $u_0=U_L$, the closed first moment $\int_Y(-u_0)\langle(K+\ddc u_0)^2\rangle$ is finite exactly when $\alpha<1/2$. Near $p$ put $t=-G_p$. Then $-u_0\asymp t^\alpha$ and $\chi''(G_p)\asymp t^{\alpha-2}$, while angular integration of $\mathrm dG_p\wedge\dc G_p$ is comparable to $\mathrm dt$. The singular contribution is thus comparable to $\int^\infty t^{2\alpha-2}\,\mathrm dt$. The smooth contributions are integrable, and the mass formula includes the divergent range.

\subsection{Ambient ideal orders and ordinary type}

To realize these data as intrinsic traces, set $X=Y\times\mathbb P^1$ and identify $Y$ with $Y\times\{z_0\}$ for a fixed $z_0$. Write $w_3$ for the third-factor form and let $\operatorname{pr}_Y\colon X\to Y$ be the projection. Assume $b<c$, and choose $L$ with $\epsilon_L<(c-b)/2$, as in the proof of \cref{thm:nonclosed-ruled-ordinary-volumes}. Put
\[
 \Theta=\operatorname{pr}_Y^*\theta+w_3,
 \qquad u_X=\operatorname{pr}_Y^*(bG_p+U_L),
 \qquad D_X=D\times\mathbb P^1.
\]
Then $u_X\in\PSH(X,\Theta)$ and $\nu(u_X,Y)=0$.

\begin{proposition}\label{prop:ruled-intrinsic-ideal-orders}\index{Multiplier ideal!unbounded ruled trace}
For every real $s>0$,
\[
 \mathcal I(su_X)=\mathcal O_X(-\lfloor sb\rfloor D_X),
 \qquad \mathcal I(su_X)\mathcal O_Y
       =\mathcal O_Y(-\lfloor sb\rfloor D).
\]
For every prime divisor $E$ on every smooth compact modification $\pi\colon Y'\to Y$, the intrinsic trace data are
\[
 \tau_{Y,u_X}(E)=b\operatorname{ord}_E(\pi^*D).
\]
Nevertheless the ordinary type of $bG_p+U_L$ has no analytic representative.
\end{proposition}
\begin{proof}
Near $D_X$, choose a transverse coordinate $z_1$. A germ whose exact transverse order is $j$ has a nonzero leading tangential coefficient. With $t=-\log|z_1|^2$, its radial integrability against $\mathrm e^{-su_X}$ is determined by
\[
 \exp\bigl(-(j+1-sb)t+s t^\alpha\bigr).
\]
Bounded errors in this exponent do not change convergence. It is integrable exactly when $j+1>sb$, and diverges at equality. Local holomorphic bounds give sufficiency, and a region where the leading tangential coefficient is nonzero gives necessity by Fubini. This proves the ambient ideal formula in the book's convention $|f|^2\mathrm e^{-su_X}$. Restriction and pullback of this principal ideal give order $\lfloor sb\rfloor\operatorname{ord}_E(\pi^*D)$. Divide by $s$ and use \cref{thm:intrinsic-trace-profile}.

The residual $U_L$ has trivial multiplier ideals at every weight, but is unbounded. Near a generic point of $D$, a potential with analytic singularities and generic Lelong number $b$ along $D$ is $b\log|z_1|^2+O(1)$. Our residual tends to $-\infty$ like $-t^\alpha$, excluding ordinary equivalence. This also covers $b=0$.
\end{proof}

\subsection{Pullback and strict representatives}

The lower volume above is positive, and its type contains a strict current on $Y$. A modification preserves both optimized masses but can destroy this last property when the ordinary type is fixed.

\Needspace{12\baselineskip}
\begin{proposition}\label{prop:ruled-pullback-ordinary-volumes}\index{Singularity volume!pure pullback}
For every modification $\pi\colon Y'\to Y$ and $v\in\PSH(Y,\theta)$ with $v\sim r$,
\[
 \underline{\operatorname{vol}}_{\pi^*\theta}(\pi^*v)
       =\underline{\operatorname{vol}}_\theta(v),\qquad
 \overline{\operatorname{vol}}_{\pi^*\theta}(\pi^*v)
       =\overline{\operatorname{vol}}_\theta(v).
\]
These masses and the ideal orders are also unchanged by a smooth gauge $(\theta,v,r)\mapsto(\theta+\ddc f,v-f,r-f)$.
\end{proposition}
\begin{proof}
By \cref{prop:psh-descent-modification}, every $\pi^*\theta$-plurisubharmonic potential descends. Since $\pi$ is surjective, ordinary inequalities are equivalent before and after pullback. Thus the competitors are in bijection. On the locus where $\pi$ is an isomorphism, locality identifies their products. Both extended products give the excluded analytic subsets zero mass, so
\[
 \pi_*\mathrm{MA}^{\mathrm{ext}}_{\pi^*\theta}(\pi^*v)
     =\mathrm{MA}^{\mathrm{ext}}_\theta(v).
\]
Taking the infimum and supremum proves the identities. Smooth gauge preserves the curvature currents and the measures; it preserves multiplier ideals because $\mathrm e^{sf}$ is locally bounded above and below.
\end{proof}

\begin{proposition}\label{prop:ruled-blowup-strict-type-obstruction}\index{Strict current!ordinary trace obstruction}\index{Blowup!centers at poles}
Suppose $b<c$ and let $\pi\colon Y'\to Y$ be a nonempty finite composition
of point blowups. There is a strict current $\pi^*\theta+\ddc v'$
with $v'\leq\pi^*r+C$, but none with $v'\sim\pi^*r$.
Both ordinary-type volumes retain the values in
\eqref{eq:nonclosed-ruled-ordinary-volumes}. Blowup centers may lie on
$D$ or on earlier exceptional curves.
\end{proposition}
\begin{proof}
Put $\beta=(c-b)w_1+aw_2$. The representatives of
\cref{lem:ruled-sublog-small-curvature} show that $U$ satisfies
\eqref{eq:small-curvature-type}; by
\cref{lem:small-curvature-zero-lelong} so does $\pi^*U$.
The logarithmic divisor is $b\pi^*D$, whose smooth curvature is
$b\pi^*w_1$, so its residual background is $\pi^*\beta$.
If a strict current existed in the ordinary type of $\pi^*r$,
\cref{thm:small-curvature-strict-type} would give a smooth $f$ with
$\pi^*\beta+\ddc f>0$. The exceptional curve $E$ of the last blowup
is contracted by the composite map. Restriction to this smooth curve
would give $\ddc(f|_E)>0$, contrary to its zero integral by Stokes.
Only the smooth function $f$ is restricted, so this argument applies
also when the original potential has a pole along $E$.

For the one-sided assertion, write $\pi=\pi_1\circ\cdots\circ\pi_N$
with point blowups $\pi_j\colon Y_j\to Y_{j-1}$ and $Y_0=Y$, and start
with the strict representative $v_0=bG_p+U_L$, $\epsilon_L<(c-b)/2$, from
the proof of \cref{thm:nonclosed-ruled-ordinary-volumes}.
Suppose $v_j$ is strict on $Y_j$, with lower bound $\delta_jh_j$ for a
Hermitian form $h_j$ and $\delta_j>0$, and let $\pi_{j+1}$ be the next
point blowup. For its exceptional divisor $E_{j+1}$, choose a smooth
metric on $\mathcal O(E_{j+1})$ whose curvature $\eta_{j+1}$ is negative
along $E_{j+1}$, which is possible since this bundle has degree $-1$ on
$E_{j+1}$, and write $\ddc\ell_{j+1}=[E_{j+1}]-\eta_{j+1}$.
For all sufficiently small $t_j>0$,
$\pi_{j+1}^*h_j-t_j\eta_{j+1}>0$.
To see this, use a blowup chart $(x,z)\mapsto(x,xz)$, in which
$E_{j+1}=\{x=0\}$ and $z$ is the vertical coordinate along $E_{j+1}$.
There the pullback has a uniformly positive horizontal coefficient, mixed coefficient
$O(|x|)$ and Schur complement at least $c_1|x|^2$. The perturbation
adds a positive vertical term at least $t_jc_2$, while its error in the
Schur complement is bounded by
$C(t_j|x|+t_j|x|^2+t_j^2)$.
Shrink the exceptional neighborhood and then choose $t_j$ small;
on its compact complement positivity persists under small perturbations.
Thus
\[
 v_{j+1}=\pi_{j+1}^*v_j+\delta_jt_j\ell_{j+1}
\]
is strict. The pullbacks are quasi-plurisubharmonic and not identically
$-\infty$, even when a center is a pole. The logarithmic terms are bounded above,
so induction gives $v'\leq\pi^*r+C$.
Finally \cref{prop:ruled-pullback-ordinary-volumes} preserves both
volumes in the pulled-back ordinary type. The additional exceptional
poles of the one-sided witness are not competitors in that fixed type.
\end{proof}

\enlargethispage{4\baselineskip}
For example take $c=1$, $b=1/4$, $\alpha=3/4$ and $a=2+\tau$, where $\tau=(|z|^2-1)/(|z|^2+1)$ is the height on the first sphere. The lower and upper volumes are $3/2$ and $9/2$. The choice $L=31$ has $\epsilon_L=(3/4)32^{-1/4}<3/8$ and gives a strict representative with infinite closed first moment. The intrinsic trace data are $\frac14\operatorname{ord}_E(\pi^*D)$. Every nonempty finite composition of point blowups, including centers on $D$, preserves both volumes and permits one-sided strictness while excluding a strict representative in the same ordinary type.
